\begin{frame}[t]{\ev{\tagO\,\tagP}Pairwise correlation does not determine best-of-$N$ coverage}
\vspace{0pt}
{\small\textbf{Calculation.} Three candidates. In each vector, 1 = candidate wrong. Listed vectors are equally likely.\par}
\vspace{.1cm}
{\small
\begin{tabularx}{\linewidth}{@{}YL{3.1cm}L{3.1cm}L{1.8cm}@{}}
\toprule
\textbf{Error vectors} & \textbf{Individual error rate} & \textbf{Pairwise correlation} & \textbf{All wrong}\\
\midrule
$\{000,\,011,\,101,\,110\}$ & 50\% & 0 & 0\%\\
all eight vectors & 50\% & 0 & 12.5\%\\
$\{001,\,010,\,100,\,111\}$ & 50\% & 0 & 25\%\\
\bottomrule
\end{tabularx}\par}
\vspace{.25cm}
{\small Same marginals and zero pairwise correlation, yet the all-wrong rate ranges from 0\% to 25\%.\\[2pt]
A selector that returns one candidate's answer succeeds with probability at most $1-P(\text{all wrong})$.\\[2pt]
An effective sample size describes the precision of an average. It is not a guarantee for search.\par}
\claim{Proposed: record every candidate's pass or fail.}
\sources{Calculation by enumeration. Ceiling: \href{https://arxiv.org/abs/2606.27288}{Chen, arXiv:2606.27288 (2026 preprint)}, ceiling $1-\beta$, $\beta$ the all-fail rate.}
\note{[0:45] Pairwise correlation does not determine best-of-N coverage. Every candidate is wrong half the time, and every pair has zero correlation. The chance that all three are wrong is zero, one in eight, or one in four. Chen's 2026 preprint shows that a selector returning one candidate cannot beat one minus the all-wrong rate. We propose recording each candidate's pass or fail to measure it. For averages, correlation does have a precise effect, which the next slide states.}
\end{frame}
